% Part: lambda-calculus
% Chapter: representability
% Section: currying

\documentclass[../../../include/open-logic-section]{subfiles}

\begin{document}

\olfileid{lam}{rep}{cur} 
\olsection{Currying（柯里化）}

$\lambd$-抽象 $\lambd[x][M]$ 表示单自变元函数，若要定义接受多个自变元的函数，这无疑是个不小的限制。一种解决办法是扩展 $\lambd$-演算，允许构成有序对、三元组等；这样一来，例如，三元函数 $\lambd[x][M]$ 便会要求其自变元是一个三元组。然而，通过\emph{柯里化}来实现这一点更为方便。

考虑一个例子。暂且假设 $\lambd$-演算中有 $+$ 运算。加法函数是二元函数，即它接受两个自变元。但 $\lambd$-抽象只提供单自变元函数：语法不允许 $\lambd[(x,y)][(x+y)]$ 这样的表达式。不过，我们可以考虑由 $\lambd[y][(x+y)]$ 给出的单自变元函数 $f_x(y)$，它将 $x$ 加到其唯一的自变元 $y$ 上。实际上，这不是单个函数，而是对应每个数 $x$ 的一族不同的“加 $x$”函数。现在我们可以定义另一个单自变元函数 $g$ 为 $\lambd[x][f_x]$。作用于自变元 $x$ 时，$g(x)$ 返回函数 $f_x$——因此其值为其他函数。若将 $g$ 作用于 $x$，再将结果作用于 $y$，便得到：$(g(x))y = f_x(y) =
x+y$。通过这种方式，单自变元函数 $g$ 就能起到与二元加法函数相同的作用。“柯里化”指的就是这种将二元函数转化为单自变元函数（其值本身也是单自变元函数）的技巧。

下面给出一个严格符合 $\lambd$-演算语法的例子。如何表示函数 $f(x,y) = x$？若要定义一个接受两个自变元并返回第一个自变元的函数，可以写 $\lambd[x][\lambd[y][x]]$，它字面上是一个接受自变元 $x$ 并返回函数 $\lambd[y][x]$ 的函数。函数 $\lambd[y][x]$ 接受另一个自变元 $y$，但将其丢弃，并总是返回 $x$。来看看将 $\lambd[x][\lambd[y][x]]$ 作用于两个自变元时会发生什么：
\begin{align*}
  (\lambd[x][\lambd[y][x]])MN 
  \bredone & (\lambd[y][M])N \\
  \bredone & M
\end{align*}

一般地，要写出由某个项 $N$ 定义的、参数为 $x_1$, \dots,~$x_n$ 的函数，可以写 $\lambd[x_1][\lambd[x_2][\ldots\lambd[x_n][N]]]$。若将 $n$ 个自变元作用于它，便得到：
\begin{multline*}
  (\lambd[x_1][\lambd[x_2][\ldots\lambd[x_n][N]]]) M_1 \dots M_n \bredone\\
  \begin{aligned}
  \bredone {} & (\Subst{(\lambd[x_2][\ldots\lambd[x_n][N]])}{M_1}{x_1}) M_2
  \dots M_n\\
   \eqs {} & (\lambd[x_2][\ldots\lambd[x_n][\Subst{N}{M_1}{x_1}]]) M_2
            \dots M_n \\
  \vdots & \\
  \bredone {} &\Subst{\Subst{P}{M_1}{x_1}\ldots}{M_n}{x_n}
  \end{aligned}
\end{multline*}
最后一行的字面意思是将函数定义体中的 $x_i$ 替换为 $M_i$，这正是将多个自变元作用于一个函数时期望的操作。
\end{document}